% from ega4-2-fr.tex
\oldpage[IV-2]{228}

\EGAIVtwoTerm{Préschéma géométriquement unibranche en un point,
géométriquement unibranche}{6.15.1}
\EGAIVtwoTerm{Préschéma intègre en un point}{3.2.4}
\EGAIVtwoTerm{Préschéma irredondant}{3.2.4}
\EGAIVtwoTerm{Préschéma localement immersible dans un schéma régulier}{5.11.1}
\EGAIVtwoTerm{Préschéma ponctuellement intègre}{4.6.9}
\EGAIVtwoTerm{Préschéma possédant la propriété $(R_k)$, possédant la propriété
$(R_k)$ en un point}{5.8.2}
\EGAIVtwoTerm{Préschéma (sur un corps) possédant la propriété $(R_k)$
géométrique en un point, possédant la propriété $(R_k)$ géométrique}{6.7.6}
\EGAIVtwoTerm{Préschéma possédant la propriété $(S_k)$, possédant la propriété
$(S_k)$ en un point}{5.7.2}
\EGAIVtwoTerm{Préschéma régulier en codimension $k$}{5.8.2}
\EGAIVtwoTerm{Préschéma (sur un corps) séparable}{4.6.2}
\EGAIVtwoTerm{Préschéma (sur un corps) séparable en un point}{4.6.9}
\EGAIVtwoTerm{Préschéma unibranche en un point, préschéma unibranche}{6.15.1}
\EGAIVtwoTerm{Préschéma universellement caténaire}{5.6.3}
\EGAIVtwoTerm{Préschéma (sur un corps) universellement réduit}{4.6.2}
\EGAIVtwoTerm{Profondeur d'un $\mathcal O_X$-Module en un point}{5.7.1}
\EGAIVtwoTerm{Propriété $P$ du premier type (du second type) associée à une
propriété $R$ d'anneaux locaux noethériens}{7.3.10}
\EGAIVtwoTerm{Propriété $P$ géométrique}{7.3.6}

\medskip
\EGAIVtwoTerm{Résolution des singularités}{7.9.1}

\medskip
\EGAIVtwoTerm{Sous-ensemble (d'un préschéma) défini sur un corps}{4.8.4}
\EGAIVtwoTerm{Sous-$\mathcal O_X$-Module défini sur un corps}{4.8.4}
\EGAIVtwoTerm{Sous-$\mathcal O_X$-Module primaire d'un $\mathcal O_X$-Module}{3.2.4}
\EGAIVtwoTerm{Sous-préschéma défini sur un corps}{4.8.4}
\EGAIVtwoTerm{Sous-préschéma fermé primaire d'un préschéma}{3.2.4}

\medskip
\EGAIVtwoTerm{$Z$-clôture d'un $\mathcal O_X$-Module}{5.9.11}

\endgroup

\clearpage
